\chapter{\ergm{} definitions reproduced on purpose}\label{app:quirks}\label{sec:quirks}

\nwergm{} is built to give the same numbers as \ergm{}. Where a definition or a computation of \ergm{}~\ergmver{} has a feature that a reader might not expect, \nwergm{} reproduces it on purpose, because a different definition would give statistics and estimates that no longer match \ergm{}'s. Where \ergm{} is plainly wrong, \nwergm{} computes the correct value and says so. This appendix lists the cases that matter to a user who compares the output of the two programs, so that a difference from the theory in the literature is not mistaken for a mistake in \nwergm{}. Each was found by comparing statistics or estimates with those of \ergm{} (Chapter~\ref{ch:validation}).
\begin{itemize}
  \item \textbf{\texttt{triangle} on a directed network.} \ergm{} accepts \texttt{triangle} on a directed network and counts the transitive plus the cyclic triples; \nwergm{} does the same, for \opt{triangle} and for \opt{structural(triangle)}. In \ergm{}'s change statistic a node that is both an in- and an out-neighbor of the head is visited twice and the change adds the ties in both directions; \nwergm{} copies this (statistic and pseudolikelihood agree with R).
  \item \textbf{Single precision.} \texttt{tripercent} computes its ratios in single precision in \ergm{}, and its observed statistic is the sum of the single-step changes along a build-up of the network; \nwergm{} copies both, and agrees with R's statistic to $2\times10^{-7}$.
  \item \textbf{Scaling by the observed network.} \texttt{degcrossprod} and \texttt{degcor} are scaled by quantities of the observed network at the moment of model creation (twice the number of ties, the variance of the observed degrees), and for \texttt{degcor} the summary statistic is the true degree correlation whereas the change statistic is an approximation. \nwergm{} copies this; the terms are therefore not available in \texttt{nwergm simulate}, where there is no observed network.
  \item \textbf{\texttt{Prod()}.} \ergm{} computes the product as \texttt{Exp(Sum(Log()))} with $\log0=-1/\sqrt{\epsilon}$. Built up tie by tie from the empty network, the statistic passes through values of about $7\times10^{7}$ and loses eight digits: for Mesa High School R gives 522.48718785 for a product whose exact value is $(3491\cdot659\cdot62)^{1/3}=522.48719458$, the value that \nwergm{} computes. The tests of \texttt{Prod()} therefore use the tolerance $10^{-6}$ for the statistics; the pseudolikelihood coefficients agree to $10^{-5}$.
  \item \textbf{\texttt{Symmetrize(rule="weak")}.} \texttt{summary()} of a formula with this operator is right, because \ergm{} adds the arcs one by one. But \texttt{ergm()} and \texttt{ergmMPLE()} start from the observed network and fill the auxiliary network with an initialization rule that, for \texttt{weak}, toggles an arc unless its reverse exists, so that every reciprocated dyad is missing from the symmetrized network on which the pseudolikelihood change statistics are computed (108 of 182 change statistics differed from brute-force differences of \texttt{summary()} on a 14-node directed network; the rules \texttt{strong}, \texttt{upper} and \texttt{lower} agree). \nwergm{} computes the exact change statistics, so that its pseudolikelihood equals R's only on a network without reciprocated ties. The test therefore compares \texttt{rule(weak)} on networks with reciprocated ties with a pseudolikelihood fitted to R's own \texttt{summary()} differences, and on a network without reciprocated ties with \texttt{ergm(estimate="MPLE")} itself. This is not to be ``fixed'' to reproduce R's initial network.
  \item \textbf{\texttt{gwesp} and \texttt{esp} with \texttt{type="RTP"}.} \ergm{}~\ergmver{} returns wrong values for the reciprocated two-path definition by default (a defect that is corrected in the development version); \nwergm{} agrees with R when R is run with \texttt{term.options=list(cache.sp=FALSE)}, to $10^{-7}$.
  \item \textbf{The penalized pseudolikelihood.} In \ergm{}~\ergmver{} \texttt{MPLE.type="penalized"} stops with an error under R~4.3 or later (a logical of length 3 in a condition); the development version fixes it. \nwergm{} reproduces the fixed behavior, which was compared with the development version to $10^{-6}$ on five networks, one of them separated.
  \item \textbf{\texttt{mcmc.diagnostics()}.} The multivariate spectral variance used for the Hotelling omnibus test and the joint Geweke test is that of \ergm{}~\ergmver{} (a Yule--Walker autoregression); the development version uses an ordinary-least-squares vector autoregression, which changes the second or third digit of these $p$-values.
  \item \textbf{Constraints.} R refuses \texttt{blockdiag} on non-contiguous blocks under the default \texttt{noncontig="merge"}, and \nwergm{} too. The constraint \texttt{hamming} is broken in R and refused; \texttt{odegreedist} on a directed network and the degree-distribution constraint on two-mode networks fail in R's released code and are refused. R's \texttt{degrees} and \texttt{degreedist} proposals for directed networks are not exactly symmetric; \nwergm{} reproduces the same chain, which agrees with R's \texttt{simulate()} within Monte Carlo error but not exactly with an exact enumeration (by at most 1\% on a 4-node directed network).
  \item \textbf{Contrastive divergence.} \ergm{} reports the covariance of a contrastive-divergence estimate as missing; \nwergm{} stores the covariance of the sample (Section~\ref{sec:cd}).
  \item \textbf{Separation.} Where a level has no node or is perfectly predicted, R reports \texttt{-Inf} or \texttt{NA} and \nwergm{} reports 0 with a standard error of 0 (for example for \texttt{idegree} and \texttt{sender}).
\end{itemize}

